% =============================================================================
%  preamble.tex —— 中译注本《无穷范畴讲义》导言区
%  底本：Lin Chen, Lectures on Infinity Categories, 2024（27 讲，203 页）
%  编译：xelatex main.tex ×3（本文件与 main.tex 由 Lead 维护，译手不得改动）
% =============================================================================
\documentclass[UTF8, fontset=fandol, 11pt, a4paper, oneside]{ctexart}

\usepackage{amsmath}
\usepackage{amssymb}
\usepackage{amsthm}
\usepackage{mathtools}
\usepackage{bm}
\usepackage{mathrsfs}
\usepackage[a4paper,left=2.6cm,right=2.6cm,top=2.8cm,bottom=2.8cm,headsep=0.5cm,footskip=1.1cm]{geometry}
\usepackage{graphicx}
\usepackage{xcolor}
\usepackage{enumitem}
\usepackage{booktabs}
\usepackage{longtable}
\usepackage{array}
\usepackage{multirow}
\usepackage{tikz}
\usepackage{tikz-cd}
\usepackage{fancyhdr}
\usepackage{newunicodechar}
% 正文中不慎以文本模式出现的符号，自动转成数学模式
\newunicodechar{∞}{\ensuremath{\infty}}
\newunicodechar{≃}{\ensuremath{\simeq}}
\newunicodechar{→}{\ensuremath{\to}}
\newunicodechar{←}{\ensuremath{\leftarrow}}
\newunicodechar{⊗}{\ensuremath{\otimes}}
\newunicodechar{⊕}{\ensuremath{\oplus}}
\newunicodechar{×}{\ensuremath{\times}}
\newunicodechar{≅}{\ensuremath{\cong}}
\newunicodechar{≠}{\ensuremath{\neq}}
\newunicodechar{≤}{\ensuremath{\leq}}
\newunicodechar{≥}{\ensuremath{\geq}}
\newunicodechar{∈}{\ensuremath{\in}}
\newunicodechar{⊂}{\ensuremath{\subset}}
\newunicodechar{∘}{\ensuremath{\circ}}
\newunicodechar{∆}{\ensuremath{\Delta}}
\newunicodechar{Λ}{\ensuremath{\Lambda}}
\newunicodechar{∂}{\ensuremath{\partial}}
\newunicodechar{π}{\ensuremath{\pi}}
\newunicodechar{α}{\ensuremath{\alpha}}\newunicodechar{β}{\ensuremath{\beta}}
\newunicodechar{γ}{\ensuremath{\gamma}}\newunicodechar{δ}{\ensuremath{\delta}}
\newunicodechar{ε}{\ensuremath{\varepsilon}}\newunicodechar{ζ}{\ensuremath{\zeta}}
\newunicodechar{η}{\ensuremath{\eta}}\newunicodechar{θ}{\ensuremath{\theta}}
\newunicodechar{ι}{\ensuremath{\iota}}\newunicodechar{κ}{\ensuremath{\kappa}}
\newunicodechar{λ}{\ensuremath{\lambda}}\newunicodechar{μ}{\ensuremath{\mu}}
\newunicodechar{ν}{\ensuremath{\nu}}\newunicodechar{ξ}{\ensuremath{\xi}}
\newunicodechar{ρ}{\ensuremath{\rho}}\newunicodechar{σ}{\ensuremath{\sigma}}
\newunicodechar{τ}{\ensuremath{\tau}}\newunicodechar{υ}{\ensuremath{\upsilon}}
\newunicodechar{φ}{\ensuremath{\varphi}}\newunicodechar{χ}{\ensuremath{\chi}}
\newunicodechar{ψ}{\ensuremath{\psi}}\newunicodechar{ω}{\ensuremath{\omega}}
\newunicodechar{Γ}{\ensuremath{\Gamma}}\newunicodechar{Θ}{\ensuremath{\Theta}}
\newunicodechar{Ξ}{\ensuremath{\Xi}}\newunicodechar{Π}{\ensuremath{\Pi}}
\newunicodechar{Σ}{\ensuremath{\Sigma}}\newunicodechar{Φ}{\ensuremath{\Phi}}
\newunicodechar{Ψ}{\ensuremath{\Psi}}\newunicodechar{Ω}{\ensuremath{\Omega}}
\usepackage[colorlinks=true,linkcolor=blue!45!black,citecolor=blue!45!black,urlcolor=blue!45!black,bookmarksnumbered=true]{hyperref}

% ---------------------------------------------------------------------------
% 版式
% ---------------------------------------------------------------------------
\linespread{1.28}
\setlength{\emergencystretch}{3em}
\allowdisplaybreaks[3]
\setlength{\parindent}{2em}
\setlength{\parskip}{0.15em}
\setcounter{tocdepth}{1}
\setcounter{secnumdepth}{3}

\pagestyle{fancy}
\fancyhf{}
\fancyhead[L]{\small\songti 无穷范畴讲义}
\fancyhead[R]{\small\songti\leftmark}
\fancyfoot[C]{\small\thepage}
\renewcommand{\headrulewidth}{0.4pt}
\renewcommand{\footrulewidth}{0pt}

\renewcommand{\contentsname}{目\hspace{2em}录}
\renewcommand{\figurename}{图}
\renewcommand{\tablename}{表}
\renewcommand{\proofname}{证明}
\renewcommand{\appendixname}{附录}
\renewcommand{\sectionmark}[1]{}
\ctexset{
  section/format = {\Large\bfseries\centering},
  section/aftername = {.\quad},
  section/number = \arabic{section},
  subsection/format = {\large\bfseries\raggedright},
}

% ---------------------------------------------------------------------------
% 讲一级标题（不编号）
% ---------------------------------------------------------------------------
\newcounter{lecture}
\newif\iflecapp
\newcommand{\lecture}[2]{%
  \clearpage
  \stepcounter{lecture}%
  \setcounter{section}{0}%
  \setcounter{nitem}{0}%
  \setcounter{equation}{0}%
  \setcounter{footnote}{0}%
  \lecappfalse
  \renewcommand{\thesection}{\arabic{section}}%
  \renewcommand{\theHsection}{lec\thelecture.\arabic{section}}%
  \section*{第 #1 讲\quad #2}%
  \addcontentsline{toc}{section}{第 #1 讲\quad #2}%
  \markboth{第 #1 讲\quad #2}{第 #1 讲\quad #2}%
}

% 附录小节：打印“附录 A　……”，并让 nitem 编号变成 A.1
\newcommand{\appendixsection}[1]{%
  \par
  \iflecapp
    \refstepcounter{section}%
  \else
    \setcounter{section}{0}%
    \renewcommand{\thesection}{\Alph{section}}%
    \lecapptrue
    \refstepcounter{section}%
  \fi
  \setcounter{nitem}{0}%
  \renewcommand{\theHsection}{lec\thelecture.\Alph{section}}%
  \section*{附录 \thesection\quad #1}%
  \addcontentsline{toc}{section}{附录 \thesection\quad #1}%
  \markright{附录 \thesection\quad #1}%
}

% ---------------------------------------------------------------------------
% 编号：每节一个统一计数器，原书所有编号对象共用“节.序号”
% ---------------------------------------------------------------------------
\newcounter{nitem}[section]
\renewcommand{\thenitem}{\thesection.\arabic{nitem}}
\renewcommand{\theHnitem}{lec\thelecture.\thesection.\arabic{nitem}}

\newtheoremstyle{cnplain}{0.9em}{0.6em}{\normalfont}{}{\bfseries}{.}{0.5em}{}
\theoremstyle{cnplain}
\newtheorem{theorem}[nitem]{定理}
\newtheorem{proposition}[nitem]{命题}
\newtheorem{lemma}[nitem]{引理}
\newtheorem{corollary}[nitem]{推论}
\newtheorem{conjecture}[nitem]{猜想}
\newtheorem{claim}[nitem]{断言}
\newtheorem{slogan}[nitem]{口号}
\newtheorem{question}[nitem]{问题}
\theoremstyle{definition}
\newtheorem{definition}[nitem]{定义}
\newtheorem{example}[nitem]{例}
\newtheorem{exercise}[nitem]{练习}
\newtheorem{challenge}[nitem]{挑战}
\newtheorem{construction}[nitem]{构造}
\newtheorem{variant}[nitem]{变体}
\newtheorem{remark}[nitem]{注}
\newtheorem{notation}[nitem]{记号}
\newtheorem{warning}[nitem]{警告}
\newtheorem{convention}[nitem]{约定}
\newtheorem{theoremdef}[nitem]{定理—定义}

% 普通编号段落：2.1. 正文……
\newcommand{\nitem}{%
  \par\refstepcounter{nitem}%
  \noindent\hspace*{2em}\thenitem.\hspace{0.5em}\ignorespaces
}

\numberwithin{equation}{section}

% 图占位（Lead 统一重绘后替换）
\newcommand{\FIGPLACEHOLDER}[1]{\fbox{\small 图 #1（待重绘）}}

% 译注：左侧竖线小字；\showyizhufalse 可隐藏
\newif\ifshowyizhu
\showyizhutrue
\newcommand{\yizhu}[1]{\ifshowyizhu\par\smallskip\noindent{\small\color{black!72}%
  \hspace*{0.4em}\vrule width 0.8pt\hspace{0.6em}\begin{minipage}[t]{0.9\linewidth}\raggedright #1\end{minipage}}\smallskip\fi}

% ---------------------------------------------------------------------------
% 常用记号（全书统一；译手只许使用这里定义的宏，见 STYLE_GUIDE.md §3）
% ---------------------------------------------------------------------------
\newcommand{\Z}{\mathbb{Z}}
\newcommand{\Q}{\mathbb{Q}}
\newcommand{\R}{\mathbb{R}}
\newcommand{\C}{\mathbb{C}}
\newcommand{\N}{\mathbb{N}}
\newcommand{\PP}{\mathbb{P}}
\newcommand{\A}{\mathbb{A}}
\newcommand{\GG}{\mathbb{G}}
\newcommand{\OO}{\mathcal{O}}
\newcommand{\cA}{\mathcal{A}}\newcommand{\cB}{\mathcal{B}}\newcommand{\cC}{\mathcal{C}}
\newcommand{\cD}{\mathcal{D}}\newcommand{\cE}{\mathcal{E}}\newcommand{\cF}{\mathcal{F}}
\newcommand{\cG}{\mathcal{G}}\newcommand{\cH}{\mathcal{H}}\newcommand{\cI}{\mathcal{I}}
\newcommand{\cJ}{\mathcal{J}}\newcommand{\cK}{\mathcal{K}}\newcommand{\cL}{\mathcal{L}}
\newcommand{\cM}{\mathcal{M}}\newcommand{\cN}{\mathcal{N}}\newcommand{\cP}{\mathcal{P}}
\newcommand{\cQ}{\mathcal{Q}}\newcommand{\cR}{\mathcal{R}}\newcommand{\cS}{\mathcal{S}}
\newcommand{\cT}{\mathcal{T}}\newcommand{\cU}{\mathcal{U}}\newcommand{\cV}{\mathcal{V}}
\newcommand{\cW}{\mathcal{W}}\newcommand{\cX}{\mathcal{X}}\newcommand{\cY}{\mathcal{Y}}
\newcommand{\cZ}{\mathcal{Z}}
\newcommand{\fA}{\mathfrak{A}}\newcommand{\fC}{\mathfrak{C}}\newcommand{\fD}{\mathfrak{D}}
\newcommand{\fM}{\mathfrak{M}}\newcommand{\fP}{\mathfrak{P}}\newcommand{\fg}{\mathfrak{g}}
\DeclareMathOperator{\Spec}{Spec}
\DeclareMathOperator{\Hom}{Hom}
\DeclareMathOperator{\Map}{Map}
\DeclareMathOperator{\Fun}{Fun}
\DeclareMathOperator{\End}{End}
\DeclareMathOperator{\Aut}{Aut}
\DeclareMathOperator{\id}{id}
\DeclareMathOperator{\coker}{coker}
\DeclareMathOperator{\im}{im}
\DeclareMathOperator{\Ker}{Ker}
\DeclareMathOperator{\Ob}{Ob}
\DeclareMathOperator{\Mor}{Mor}
\DeclareMathOperator{\colim}{colim}
\DeclareMathOperator{\hocolim}{hocolim}
\DeclareMathOperator{\holim}{holim}
\DeclareMathOperator{\Tot}{Tot}
\DeclareMathOperator{\Ind}{Ind}
\DeclareMathOperator{\Pro}{Pro}
\DeclareMathOperator{\PSh}{PSh}
\DeclareMathOperator{\Sh}{Sh}
\DeclareMathOperator{\Sing}{Sing}
\DeclareMathOperator{\Sp}{Sp}
\DeclareMathOperator{\Spc}{Spc}
\DeclareMathOperator{\Cat}{Cat}
\DeclareMathOperator{\Set}{Set}
\DeclareMathOperator{\sSet}{sSet}
\DeclareMathOperator{\Grpd}{Grpd}
\DeclareMathOperator{\Kan}{Kan}
\DeclareMathOperator{\Ch}{Ch}
\DeclareMathOperator{\Mod}{Mod}
\DeclareMathOperator{\Alg}{Alg}
\DeclareMathOperator{\coAlg}{coAlg}
\DeclareMathOperator{\Vect}{Vect}
\DeclareMathOperator{\Ab}{Ab}
\DeclareMathOperator{\Grp}{Grp}
\DeclareMathOperator{\Sptr}{Sptr}
\DeclareMathOperator{\Mon}{Mon}
\DeclareMathOperator{\PrSt}{PrSt}
\DeclareMathOperator{\mult}{mult}
\DeclareMathOperator{\ind}{ind}
\DeclareMathOperator{\res}{res}
\DeclareMathOperator{\oblv}{oblv}
\DeclareMathOperator{\LMod}{LMod}
\DeclareMathOperator{\Top}{Top}
\DeclareMathOperator{\Ho}{Ho}
\DeclareMathOperator{\hTop}{hTop}
\DeclareMathOperator{\hGrpd}{hGrpd}
\DeclareMathOperator{\Ring}{Ring}
\DeclareMathOperator{\ho}{ho}
\newcommand{\op}{\mathrm{op}}
\newcommand{\ot}{\otimes}
\newcommand{\we}{\simeq}
\newcommand{\ra}{\rightarrow}
\newcommand{\la}{\leftarrow}
\newcommand{\To}{\longrightarrow}
\newcommand{\from}{\longleftarrow}
\newcommand{\into}{\hookrightarrow}
\newcommand{\onto}{\twoheadrightarrow}
\newcommand{\defeq}{\coloneqq}
\newcommand{\ie}{i.e.\ }
\newcommand{\eg}{e.g.\ }
